\documentclass[11pt]{amsart}

\usepackage{amsmath, amssymb, amsthm}
\usepackage{mathtools}
\usepackage{thmtools}
\usepackage{enumitem}
\usepackage{tikz-cd}
\usepackage{hyperref}
\usepackage[nameinlink]{cleveref}

% % Theorem environments
\declaretheorem[numberwithin=section]{theorem}
\declaretheorem[style=plain, sibling=theorem]{lemma, proposition, corollary}
\declaretheorem[style=definition]{definition, example, condition}
\declaretheorem[style=remark]{remark, note, observation, todo}


\title{April 16, 2026}


\begin{document}
\maketitle

We observe the follow at dawn.
\begin{itemize}
    \item Chern--Simons invariants of connections on a closed manifold is a gauge-invariant (up to $2 \pi i \mathbb{Z}$).
    \item Chern--Simons invariants of connections on manifolds with boundary are not gauge invariants.
    The additional information we need is log-decoration with a choice of invariant line of the peripheral subgroup.
\end{itemize}
The above shows that, we do not consider the Chern--Simons invariant of a knot as an invariant of $\mathrm{SL}_2$-colored diagrams.
It should be regarded as an invariant of log-decorated diagrams, which is equivalent to the choice of log of the eigenvalue of the meridian.

Additionally, I found that I have gotten confused with the partition function and the Chern--Simons invariant:
\begin{itemize}
    \item The Chern--Simons invariant $\mathrm{CS}(A)$ is a functional of a connection $A$ on a manifold $M$.
    Although it is not invariant under the action of the gauge group in general, it is invariant for flat connections (up to $2 \pi i \mathbb{Z}$).

    \item On the other hand, the partition function is defined by the path integral
    \[
        Z(M) = \int \mathcal{D}A \, \exp \left( ik \mathrm{CS}(A) \right)
    \]
    where $k$ is the level, and the integral is computed along the space of connections on $M$.

    It is known that the expectation value of the Willson loop observable $W_R(K)$ can be computed as
    \[
        \frac{1}{Z(S^3)} \int \mathcal{D}A \exp \left( ik \mathrm{CS}(A) \right) W_R(K)
    \]
    where $R$ is a $\mathrm{SU}(2)$-representation.

    By Witten, we have
    \[
        \langle W_R(K) \rangle = V_K(q)
    \]
    where $R$ is the natural representation (of dim 2) and $V_K(q)$ is the Jones polynomial with
    \[
        q = \exp \left( \frac{2\pi i}{k + 2} \right).
    \]
    It seems like that if we replace $R$ into general spin-$j$ representations, then the expected value would be proportional to the colored Jones poklynomials though it is not confirmed.
\end{itemize}

After sleeping, my daily workout, having launch, I checked the yesterday's review file.
It was still the problem that I did not use idiomatic verb patterns.
In my guess, I should find a way to get be accustomed with idiom.
The problem using phrases sound translated might be solved by the same way.

Another noticable feedback was to keep the interpretation modest.
One of my insight is that recording impressions and intuitions on things that I have considered in each day might be valuable.
Those are good guidlines for mathematical research as propostitions we've found, but they are fragile so that they can be lost at times.

\par\vspace{1em}

I also checcked updated papers from arXiv today.

\par\vspace{1em}

Due to preparation for the talk tomorrow, I am not able to record research today.
The observations obtained by reading McPhail-Snyder's paper will be recorded later.







\bibliographystyle{amsalpha}
% \nocite{*}
\bibliography{bibliography}
\end{document}
